\documentclass{beamer}
\usepackage{ctex}
\usepackage{tikz}
\usepackage{tikz-cd}

\usepackage{unicode-math}
%\setmathfont{STIX Two Math}
\setmathfont{TeX Gyre Termes Math}
%\setmathfont{Libertinus Math}
%\setmathfont{Fira Math}

%%%%%%%% FONT TIMESAVES %%%%%%%
\newcommand{\bT}{\mathbb{T}}
\newcommand{\bN}{\mathbb{N}}
\newcommand{\bI}{\mathbb{I}}

\newcommand{\cC}{\mathcal{C}}

%%%%%%%% CATEGORY THEORY %%%%%%%
\newcommand{\id}{\mathrm{id}}
\DeclareMathOperator{\Hom}{Hom}

\newcommand{\cat}[1]{\mathrm{#1}}
\newcommand{\Set}{\cat{Set}}

\DeclareFontFamily{U}{dmjhira}{}
\DeclareFontShape{U}{dmjhira}{m}{n}{ <-> dmjhira }{}

\DeclareRobustCommand{\yo}{\text{\usefont{U}{dmjhira}{m}{n}\symbol{"48}}}

\DeclareMathOperator*{\colim}{colim}

\newcommand{\inl}{\operatorname{inl}}
\newcommand{\inr}{\operatorname{inr}}

%%%%%%%%% GEOMTERY %%%%%%%%%%%%%%
\newcommand{\Spec}{\operatorname{Spec}}
\newcommand{\Zar}{\mathrm{Zar}}

%%%%%%%%% TYPE THEORY %%%%%%%%%%%
\newcommand{\OO}{\bigcirc}
\newcommand{\UU}{\mathcal{U}}
\newcommand{\isEquiv}{\bibopenparen{isEquiv}}
\newcommand{\hlevel}{\mathrm{hlevel}}
\newcommand{\isContr}{\operatorname{isContr}}
\newcommand{\isProp}{\operatorname{isProp}}
\newcommand{\Prop}{\mathrm{Prop}}
\newcommand{\refl}{\mathrm{refl}}
\newcommand{\isModal}{\operatorname{isModal}}
\newcommand{\ap}{\operatorname{ap}}
\newcommand{\fib}{\operatorname{fib}}
\newcommand{\im}{\operatorname{im}}
\newcommand{\tot}{\mathrm{tot}}
\newcommand{\pt}{\mathrm{pt}}

\title{半单纯集(Semi-simplicial set)}

\date{}

\begin{document}

\begin{frame}
    % Print the title page as the first slide
    \titlepage
\end{frame}

\begin{frame}
对 $n:\bN$, 记
\[
[n] := \{ k : \bN \mid k < n \}
\]
\end{frame}

\begin{frame}{半单纯范畴}
半单纯范畴 $\Delta_{+}$, 其对象为自然数, 给定 $n, m : \Delta_{+}$, $n$ 与 $m$ 之间的态射是从 $[n]$ 到 $[m]$ 的严格单调递增函数.
\end{frame}

\begin{frame}{直范畴(direct category)}

$ω$ 是自然数偏序 $(\mathbb{N},≤)$ 对应的范畴

\vspace{6mm}

称一个范畴 $C$ 是直的(direct), 如果存在一个秩函子(rank functor) $\varphi : C \to \omega$, 它反射恒等态射(reflects identities), 即: 对任意 $x,y:C$, 若存在态射 $f : x \to y$ 使得 $\varphi(x) = \varphi(y)$, 则 $f$ 必为恒等态射, 即 $(y,f) = (x,id_x)$. 对于 $x:C$, 我们称 $\varphi(x)$ 为 $x$ 的秩(rank).

\vspace{6mm}

对于 $n : \mathbb{N}$，记 $C^{<n}$ 为由所有秩小于 $n$ 的对象构成的满子范畴；记 $C^{=n}$ 为秩恰好等于 $n$ 的对象的类型。它同样是 $C$ 的一个满子范畴，但由于它同时也是一个离散范畴，其范畴结构并不特别有趣。

\vspace{6mm}

几何上看，Direct Category（直范畴）可以理解成一个“只能向上生长的有向分层图（layered directed graph）”。
\end{frame}

\end{document}